\section{Conclusion}

The primary objective of this project was to explore backend architectural design choices, 
including system architecture, database strategies, and caching mechanisms, 
through both theoretical research and practical implementation. 
By designing and developing a cuisine-based application, 
the project provided a concrete context in which these concepts could be examined and applied.

Through researching architectural patterns and backend optimization techniques, 
I developed a stronger understanding of the trade-offs involved in system design. 
The process of implementing the software strengthened this understanding, 
since it involved translating theoretical principles into practical technical solutions.
This combination of theoretical study and hands-on implementation 
was essential in achieving the project's learning objectives.

In addition, the project highlighted the importance of performance considerations and system scalability through, 
data management strategies and caching methods.
Although not all techniques explored in theory were implemented, 
the process of analyzing potential alternatives provided valuable insight 
into how different design decisions can influence a software application.

However, there were also a few struggles throughout the independent study project. One of the biggest challenges was learning Java Spring Boot.
I had previous experience with Java through past courses, but learning the Spring Boot framework was quite difficult.
It differed from the Python frameworks I was used to, particularly in its emphasis on object-oriented programming,
different approaches to dependency management, and understanding the bean lifecycle.
Learning these new concepts took time before I could begin implementation, and applying them within an unfamiliar architecture made the process even more challenging.
This project was a valuable experience to explore a new language and framework, which helped me develop a more solid grasp of Java Spring Boot that I can utilize in the future.
However, if I had known that I would not be working in Korea after graduation, I would have implemented this project using a framework I am more comfortable with,
so that I could focus more on applying theoretical knowledge to code. I believe that I could have experimented with various software architectures
and caching strategies if I had initially started with a Python framework.

Although there were some aspects of the project that could have been approached differently, I am still satisfied with the results I achieved.
I was able to learn Java Spring Boot, and the project also successfully met its goal of deepening my understanding 
of backend architecture, database design, and caching strategies, 
while also demonstrating the practical challenges associated with building a complete system.
Overall, the independent study project was a valuable experience that enhanced my understanding of backend web development.
